\documentclass[../../main2.tex]{subfiles}

\begin{document}

	% ======================================================================
	% 骨架文件：小节标题与追问链为终版；正文由后续会话填充。
	% 原书材料接入点：
	%   - 原书 conclusion（循环的闭合）→ 本章，升级为「预测 → 新动机」闭环
	% ======================================================================

	\chapter{循环的闭合：未来塑造新的动机}
	\label{ch:closure}

	\begin{motivation}
		上一章遗留的问题：预测的方法与边界都交代了。但预测不是旁观——一份被足够多人相信的预测，会改变被预测的世界本身。我们预测出的未来，会不会反过来塑造新的动机？\\
		\textbf{核心动机}：闭合全书——预测的结果回流为下一轮的动机与目标（P5 $\to$ P1 的循环边）；并给出一句诚实的总结：这本书给出的是结构，不是预言。
	\end{motivation}

	\begin{logicmap}
	\begin{tikzpicture}[node distance=1.0cm, every node/.style={draw, rectangle, rounded corners, align=center}]
	\node (q) {预测会改变世界吗？};
	\node (s) [below=of q] {自指与预期（§14.1）};
	\node (b) [below=of s] {回到 Part I（§14.2）};
	\node (e) [below=of b] {结语：结构，不是预言（§14.3）};
	\node (l) [left=1.8cm of b] {Part I\\动机与目标};
	\draw[-{Stealth}] (q) -- (s);
	\draw[-{Stealth}] (s) -- (b);
	\draw[-{Stealth}] (b) -- (e);
	\draw[-{Stealth}, dashed] (b) -- (l);
	\end{tikzpicture}
	\end{logicmap}

	% ==================================================================
	\section{自指与预期}
	\label{sec:cl-selfreference}
	预期本身改变结果。这不是哲学玄想，是机制性的：当足够多的人相信某件事会发生，他们的行为就会让这件事发生（自证预言）；或者，当足够多的人相信某件事不会发生，他们的行为就会阻止它发生（自否预言）。

	银行挤兑是最经典的自证案例：如果储户相信银行会倒闭，他们就会去取钱，取钱本身导致银行倒闭——信念通过行为实现了自己。反过来，央行宣布「存款有保障」如果足够可信，就能触发自否：储户相信不会挤兑，所以不取钱，所以确实没有挤兑。两种方向的不对称性在于：恐慌的传播速度远快于信任的建立速度（坏消息走模因通道的传播强度系统性高于好消息），所以自证预言比自否预言更容易实现。

	这给第13章的方法加上了第六种失败模式：**预测改变初值。** 如果你的预测被公开，且受众的行为受该预测驱动，那么预测的发布本身就改变了系统状态——你算的是「没有这条预测时」的未来，但实际运行的是「有这条预测时」的系统。缓解手段：延迟发布、模糊化表述、或只向内部决策者披露——但都不能消除，只能降低反身性强度。

	\begin{readingnote}
		「我看空这家银行」这条推特本身就能让银行出事——如果足够多的人看到并相信。这不是造谣，是反身性：预测和行为之间的因果箭头不是单向的，而是成环的。社交媒体把这种反身性放大到了前所未有的尺度：一条推文可以在几小时内触达数亿人，而传统媒体需要几周。
	\end{readingnote}

	\paragraph*{逻辑衔接}
	预测改变动机、动机改变行为、行为改变结果——这条回流的终点在哪？下一节把循环焊上。
	% ==================================================================
	\section{回到 Part I}
	\label{sec:cl-return}
	第13章的预测结果成为新一轮的动机输入：P5 → P1，闭环成立。

	完整的一圈：动机（为什么行动）→ 目标（要什么）→ 状态（现在在哪）→ 概率（未来可能怎样）→ 贡献（哪些通道在驱动）→ 通道（具体机制）→ 预测（综合判断）→ 新动机（基于预测调整行动）。每个箭头都有独立的定义与度量（第1-13章各占一段），循环只是把它们首尾相接。

	循环有终点吗？有——第5章可预测窗口的两端。完全冻结（所有参数不变，系统不再产生新信息）和完全混沌（噪声主导，任何预测都不优于随机猜测）都不再产生「新动机」。现实系统在这两端之间运行，所以循环持续运转。

	循环不是圆圈而是螺旋：每圈的状态空间都被上一圈改写。人做出预测 → 预测改变行为 → 行为改变世界 → 新的世界产生新的动机。这就是原书说的「人化自然」：人不是被动地观察自然规律，而是主动地用预测和干预重塑自然。每转一圈，地形都变了。

	\begin{questionbox}
		「循环解释一切，等于什么都没解释？」不等于。循环不是解释，是结构声明——它告诉你「有哪些要素、怎么连接」。每个箭头上的独立定义与度量（条件概率公式、贡献分解算法、通道判据）才是真正的解释内容。循环只是说「这些解释不是孤立的，它们首尾相接形成一个闭合系统」。
	\end{questionbox}

	\paragraph*{逻辑衔接}
	循环闭合了。最后一节，把全书压成一句诚实的话。
	% ==================================================================
	\section{结语：结构，不是预言}
	\label{sec:cl-coda}
	这本书给出的是结构，不是预言。

	它能告诉你：哪几个因素在起作用（四通道）、各占多少（贡献分解）、接下来最可能怎样（条件概率链）、哪里会算不动（六种失败模式）。它不能告诉你：明天会发生什么。任何声称能告诉你「明天」的书，要么在骗你，要么在卖你。

	回放引言的检验标准：(1) 一句话立论——社会行为可以分解为有限条通道的贡献之和，每条通道有独立判据和定量方法。✓ (2) 可证伪——如果某个社会现象无法归入任何通道，或者通道贡献的测量结果与预测系统性偏差，框架失效。✓ (3) 可操作——贡献表 + 概率链给出了从数据到判断的完整流程，不需要天才直觉。✓

	我们能说什么：结构、趋势、边界。我们不能说什么：精确时间点、精确幅度、黑天鹅事件。

	最后一页，回到第一章那个问「为什么」的小孩。他问的不是「明天股市涨还是跌」，他问的是「为什么社会是这样运作的」。这本书的回答是：因为它由几条通道组成，每条通道有自己的规则，规则之间有嵌套关系，整个系统在一个不断改写的循环里运转。这不是一个让人安心的答案——它意味着没有终极原因，只有当前这一圈的结构。但它是诚实的。
\end{document}
